% Las figuras convertidas desde WMF no llevan ancho explícito. Algunas
% (diagramas grandes) superan el alto de página e invaden la cabecera;
% otras (fórmulas cortas inline) son pequeñas y deben conservar su tamaño
% natural. adjustbox "max width/height" solo reduce si hace falta, nunca
% amplía, a diferencia de fijar `width=` por defecto con \pkg{graphicx}.
\usepackage[export]{adjustbox}
\adjustboxset{max width=0.9\textwidth, max height=0.72\textheight}

% Las figuras con caption ("Figura X-Y: ...") ahora son entornos `figure`
% flotantes reales (para que la página se corte antes de la figura en
% vez de solaparla con la cabecera/texto). `float`+`H` las fija justo
% donde se referencian, sin dejarlas "flotar" lejos de su texto.
\usepackage{float}
\floatplacement{figure}{H}

% Nuestro propio texto ya incluye "Figura X-Y: ..." como pie de foto, así
% que se anula el "Figura N:" que LaTeX añadiría automáticamente para no
% duplicar la numeración.
\usepackage{caption}
\captionsetup{labelformat=empty, labelsep=none, font=it}
